\documentclass[aspectratio=169,11pt]{beamer}

% ============================================================
% PACOTES
% ============================================================

\usepackage[utf8]{inputenc}
\usepackage[T1]{fontenc}
\usepackage[brazil]{babel}

\usepackage{amsmath}
\usepackage{amssymb}
\usepackage{graphicx}
\usepackage{tikz}
\usepackage[most]{tcolorbox}
\usepackage{lmodern}

\usetikzlibrary{
    shadows,
    shadows.blur
}

% ============================================================
% CORES
% ============================================================

\definecolor{mainblue}{HTML}{3B35B9}
\definecolor{darkblue}{HTML}{16175A}

\definecolor{lightbox}{HTML}{ECECF7}

\definecolor{alertred}{HTML}{D60000}
\definecolor{alertbg}{HTML}{F8E5E5}

\definecolor{softgray}{HTML}{F7F7F9}


% ============================================================
% CONFIGURAÇÃO GERAL DO BEAMER
% ============================================================

\setbeamertemplate{navigation symbols}{}

\setbeamercolor{background canvas}{
    bg=white
}

\setbeamercolor{normal text}{
    fg=black
}

\setbeamercolor{item}{
    fg=mainblue
}

\setbeamerfont{frametitle}{
    size=\Large,
    series=\mdseries
}

\setbeamerfont{normal text}{
    size=\normalsize
}

% Margens laterais
\setbeamersize{
    text margin left=0.60cm,
    text margin right=0.60cm
}

% Marcadores
\setbeamertemplate{itemize item}{
    \color{mainblue}$\blacktriangleright$
}

\setbeamertemplate{itemize subitem}{
    \color{mainblue}$\blacktriangleright$
}

% Espaçamento das listas
\setlength{\leftmargini}{1.5em}
\setlength{\leftmarginii}{1.3em}


% ============================================================
% FAIXA SUPERIOR / TÍTULO DO SLIDE
% ============================================================

\setbeamercolor{frametitle}{
    bg=mainblue,
    fg=white
}

\setbeamertemplate{frametitle}{%
    \nointerlineskip%
    \begin{beamercolorbox}[
        wd=\paperwidth,
        ht=0.72cm,
        dp=0.17cm,
        leftskip=0.32cm
    ]{frametitle}%
        \usebeamerfont{frametitle}%
        \insertframetitle%
    \end{beamercolorbox}%
}


% ============================================================
% RODAPÉ
% ============================================================

\setbeamercolor{footerleft}{
    bg=darkblue,
    fg=white
}

\setbeamercolor{footerright}{
    bg=mainblue,
    fg=white
}

% ------------------------------------------------------------
% Estes parâmetros podem ser ajustados facilmente
% ------------------------------------------------------------

% Espaço branco mínimo entre conteúdo e rodapé
\newcommand{\footersafespace}{0.30cm}

% Altura acima da linha de base do texto
\newcommand{\footerheight}{0.25cm}

% Profundidade abaixo da linha de base.
% Aumentar este valor faz o texto subir dentro da faixa.
\newcommand{\footerdepth}{0.09cm}


\setbeamertemplate{footline}{%
    \leavevmode%
    \vbox{%
        % Área de segurança:
        % impede caixas e textos de encostarem no rodapé
        \vskip\footersafespace%
        %
        \hbox to \paperwidth{%
            %
            % Parte esquerda
            \begin{beamercolorbox}[
                wd=0.85\paperwidth,
                ht=\footerheight,
                dp=\footerdepth,
                leftskip=0.10cm
            ]{footerleft}%
                \tiny
                INF01008 -- Aula 1%
            \end{beamercolorbox}%
            %
            % Parte direita
            \begin{beamercolorbox}[
                wd=0.15\paperwidth,
                ht=\footerheight,
                dp=\footerdepth,
                rightskip=0.10cm
            ]{footerright}%
                \hfill
                \tiny
                \insertframenumber/\inserttotalframenumber%
            \end{beamercolorbox}%
        }%
    }%
}


% ============================================================
% CAIXAS
% ============================================================

\tcbset{
    slidebox/.style={
        enhanced,
        boxrule=0pt,
        arc=2mm,
        outer arc=2mm,
        left=2.2mm,
        right=2.2mm,
        top=1.4mm,
        bottom=1.4mm,
        toptitle=1mm,
        bottomtitle=1mm,
        drop shadow={
            shadow xshift=1.4mm,
            shadow yshift=-1.4mm,
            opacity=0.35
        }
    }
}


% ------------------------------------------------------------
% Caixa azul: pergunta / destaque
% ------------------------------------------------------------

\newtcolorbox{questionbox}[1]{
    slidebox,
    colback=lightbox,
    colframe=lightbox,
    coltitle=white,
    colbacktitle=darkblue,
    title=#1,
    fonttitle=\normalsize
}


% ------------------------------------------------------------
% Caixa vermelha: impacto / observação / conclusão
% ------------------------------------------------------------

\newtcolorbox{impactbox}[1]{
    slidebox,
    colback=alertbg,
    colframe=alertbg,
    coltitle=white,
    colbacktitle=alertred,
    title=#1,
    fonttitle=\normalsize
}


% ============================================================
% COMANDO PARA IMAGENS EM SLIDES
% ============================================================

% Uso:
%
% \slideimage{arquivo.pdf}
%
% ou
%
% \slideimage[width=0.85\linewidth]{arquivo.pdf}
%

\newcommand{\slideimage}[2][]{%
    \begin{center}
        \includegraphics[
            width=\linewidth,
            height=0.56\textheight,
            keepaspectratio,
            #1
        ]{#2}
    \end{center}
}


% ============================================================
% SLIDE DIVISOR DE SEÇÃO
% ============================================================

\newcommand{\sectionframe}[1]{%
    {%
    \setbeamertemplate{frametitle}{}%
    \begin{frame}
        \centering
        \vfill

        {\Large #1}

        \vfill
    \end{frame}
    }%
}


% ============================================================
% METADADOS
% ============================================================

\title{Programação Paralela}
\subtitle{Aula 1: Motivação para Paralelismo}

\author{Arthur Francisco Lorenzon}

\institute{
    Instituto de Informática -- UFRGS
}

\date{2026/2}


% ============================================================
% DOCUMENTO
% ============================================================

\begin{document}


% ============================================================
% CAPA
% ============================================================

{%
\setbeamertemplate{frametitle}{}%

\begin{frame}

    \vfill

    \begin{center}

        \begin{tcolorbox}[
            enhanced,
            width=0.92\textwidth,
            colback=mainblue,
            colframe=mainblue,
            arc=2mm,
            boxrule=0pt,
            left=3mm,
            right=3mm,
            top=5mm,
            bottom=5mm,
            drop shadow={
                shadow xshift=1.7mm,
                shadow yshift=-1.7mm,
                opacity=0.35
            }
        ]

            \centering

            {\color{white}\LARGE
                \inserttitle
                \par
            }

            \vspace{0.15cm}

            {\color{white}\Large
                \insertsubtitle
            }

        \end{tcolorbox}

        \vspace{0.75cm}

        {\large
            \insertauthor
        }

        \vspace{0.35cm}

        {\small
            \insertinstitute
        }

        \vspace{0.35cm}

        {\small
            \insertdate
        }

    \end{center}

    \vfill

\end{frame}
}


% ============================================================
% EXEMPLO: SLIDE DIVISOR
% ============================================================

\sectionframe{Por que paralelismo?}


% ============================================================
% EXEMPLO: SLIDE COM PERGUNTA
% ============================================================

\begin{frame}{O Problema Fundamental}

    \vspace{0.25cm}

    \textbf{Pergunta principal:}

    \vspace{0.35cm}

    \begin{questionbox}{Pergunta}
        Por que não simplesmente aumentar a frequência do
        processador?
    \end{questionbox}

\end{frame}


% ============================================================
% EXEMPLO: SLIDE COM EQUAÇÃO E CAIXA DE IMPACTO
% ============================================================

\begin{frame}{Limite Térmico}

    \textbf{Problema:}
    aumentar a frequência aumenta a potência.

    \vspace{0.15cm}

    \[
        P \propto C \cdot V^2 \cdot f
    \]

    Onde:

    \begin{itemize}
        \item $V$ = tensão, $f$ = frequência
        \item $C$ = capacitância
    \end{itemize}

    \vspace{0.10cm}

    \textbf{Consequência:}

    \begin{itemize}
        \item Mais clock $\rightarrow$ mais calor
        \item Resfriamento vira limitador físico
        \item Surge o conceito de TDP
    \end{itemize}

    \vspace{0.10cm}

    \begin{impactbox}{Impacto}
        Não podemos continuar aumentando a frequência
        indefinidamente.
    \end{impactbox}

\end{frame}

% ============================================================
% EXEMPLO: SLIDE DE LISTA
% ============================================================

\begin{frame}{Consumo de Energia}

    Supercomputadores modernos consomem:

    \begin{itemize}
        \item Dezenas de megawatts
        \item Custo operacional elevado
        \item Impacto ambiental significativo
    \end{itemize}

    \vspace{0.15cm}

    Exemplo:

    \begin{itemize}
        \item $1\,\mathrm{MW} \approx 1000$ casas
    \end{itemize}

    \vspace{0.15cm}

    \textbf{Problema:}
    aumentar a frequência aumenta o consumo total.

    \vspace{0.30cm}

    \begin{questionbox}{Pergunta}
        Vale mais dobrar o clock ou dobrar os núcleos
        mantendo a eficiência?
    \end{questionbox}

\end{frame}


% ============================================================
% EXEMPLO: TEXTO + IMAGEM
% ============================================================

\begin{frame}{Fim da Escalabilidade de Dennard}

    \begin{columns}[T,onlytextwidth]

        % ----------------------------------------------------
        % COLUNA ESQUERDA
        % ----------------------------------------------------

        \begin{column}{0.48\textwidth}

            Dennard (1974):

            \begin{itemize}
                \item Transistores menores
                \item Mesmo consumo por área
                \item Frequência podia aumentar
            \end{itemize}

            \vspace{0.15cm}

            Após $\approx 2005$:

            \begin{itemize}
                \item Tensão não reduz proporcionalmente
                \item Densidade de potência aumenta
                \item Frequência estagna
            \end{itemize}

        \end{column}


        % ----------------------------------------------------
        % COLUNA DIREITA
        % ----------------------------------------------------

        \begin{column}{0.47\textwidth}

            \centering

            \vspace{0.20cm}

            % Troque pelo caminho da sua imagem.
            %
            % Exemplo:
            %
            % \slideimage{figuras/dennard.pdf}

            %\fbox{%
            %    \parbox[c][0.50\textheight][c]
            %    {0.90\linewidth}{%
            %        \centering
            %        \textit{Imagem / gráfico}
            %    }%
            %}

            \slideimage[
                  width=0.93\linewidth
              ]{imagens/wallpaper.jpeg}

        \end{column}

    \end{columns}

\end{frame}


% ============================================================
% MESMO SLIDE USANDO UMA IMAGEM REAL
% ============================================================

% Descomente quando possuir o arquivo:
%
% \begin{frame}{Fim da Escalabilidade de Dennard}
%
%     \begin{columns}[T,onlytextwidth]
%
%         \begin{column}{0.48\textwidth}
%
%             Dennard (1974):
%
%             \begin{itemize}
%                 \item Transistores menores
%                 \item Mesmo consumo por área
%                 \item Frequência podia aumentar
%             \end{itemize}
%
%             Após $\approx 2005$:
%
%             \begin{itemize}
%                 \item Tensão não reduz proporcionalmente
%                 \item Densidade de potência aumenta
%                 \item Frequência estagna
%             \end{itemize}
%
%         \end{column}
%
%         \begin{column}{0.47\textwidth}
%
%             \slideimage[
%                 width=0.93\linewidth
%             ]{figuras/dennard.pdf}
%
%         \end{column}
%
%     \end{columns}
%
% \end{frame}


% ============================================================
% EXEMPLO: RESULTADO / CONCLUSÃO
% ============================================================

\begin{frame}{Resumo dos Limites Físicos}

    \begin{itemize}
        \item Limite térmico
        \item Consumo energético
        \item Fim de Dennard
        \item Memory wall
    \end{itemize}

    \vspace{0.30cm}

    \begin{impactbox}{Conclusão}
        A solução não foi aumentar o clock.

        A solução foi o paralelismo.
    \end{impactbox}

    \vspace{0.25cm}

    \begin{questionbox}{Pergunta Final}
        Paralelismo resolve todos esses problemas?
    \end{questionbox}

\end{frame}


% ============================================================
% EXEMPLO: SLIDE MATEMÁTICO
% ============================================================

\begin{frame}{Lei de Amdahl (1967)}

    Dividimos o programa em:

    \begin{itemize}
        \item Fração serial: $S$
        \item Fração paralelizável: $P$
    \end{itemize}

    Com

    \[
        S + P = 1
    \]

    Speedup com $N$ processadores:

    \[
        \operatorname{Speedup}(N)
        =
        \frac{1}{
            S + \dfrac{P}{N}
        }
    \]

\end{frame}


% ============================================================
% EXEMPLO: SLIDE DE EXEMPLO NUMÉRICO
% ============================================================

\begin{frame}{Exemplo 1}

    Suponha:

    \begin{itemize}
        \item $90\%$ paralelizável ($P=0.9$)
        \item $10\%$ serial ($S=0.1$)
        \item $N=10$
    \end{itemize}

    \vspace{0.30cm}

    \[
        \operatorname{Speedup}
        =
        \frac{1}{
            0.1 + \dfrac{0.9}{10}
        }
        =
        \frac{1}{0.19}
        \approx
        5.26
    \]

    \vspace{0.25cm}

    \begin{impactbox}{Observação}
        Mesmo com 10 processadores, não chegamos nem perto
        de $10\times$.
    \end{impactbox}

\end{frame}


\end{document}